\cleardoublepage
\phantomsection
\begin{planbox}[แผนบริหารการสอนประจำบทที่ 6\quad อุณหพลศาสตร์ ความดันไอ และหม้อนึ่งความดันไอ (สัปดาห์ที่ 6)]
\noindent\textbf{1. หัวข้อเนื้อหาประจำบท (บรรยาย 2 ชั่วโมง และปฏิบัติการ 3 ชั่วโมง)}
\begin{itemize}
    \item 6.1 กฎของแก๊สอุดมคติและความสัมพันธ์ระหว่างความดันกับอุณหภูมิในภาชนะปิด
    \item 6.2 ความดันไออิ่มตัวและการเพิ่มขึ้นของจุดเดือดของน้ำเมื่อความดันสูงขึ้น
    \item 6.3 หลักการทางอุณหพลศาสตร์ของหม้อนึ่งความดันไอ (Autoclave) ที่สภาวะมาตรฐาน $121\,^\circ\text{C}$, $15\text{ psi}$
    \item 6.4 ปฏิบัติการที่ 6 การศึกษาความสัมพันธ์ระหว่างความดันกับจุดเดือด และกรณีศึกษาการตรวจสอบรอบการทำงานของเครื่อง Autoclave
\end{itemize}

\noindent\textbf{2. วัตถุประสงค์เชิงพฤติกรรมและผลลัพธ์การเรียนรู้รายวิชา (CLOs ที่เชื่อมโยง PLOs)}
\begin{enumerate}
    \item อธิบายความสัมพันธ์ระหว่างความดันสัมบูรณ์กับอุณหภูมิเคลวินตามกฎของเกย์-ลูสแซก ($P_1 / T_1 = P_2 / T_2$) ได้ \textbf{(CLO1 เชื่อมโยง PLO2)}
    \item วิเคราะห์กลไกทางฟิสิกส์ของเครื่อง Autoclave ได้ว่าทำไมต้องไล่อากาศออกจากหม้อนึ่งให้หมด และทำไมต้องใช้ไอน้ำอิ่มตัวความดันเกจ $15\text{ psi}$ เพื่อให้ได้อุณหภูมิ $121\,^\circ\text{C}$ สำหรับทำลายเอนโดสปอร์ของแบคทีเรีย \textbf{(CLO1, CLO2 เชื่อมโยง PLO2, PLO3)}
    \item ปฏิบัติตามขั้นตอนความปลอดภัยของภาชนะรับความดันสูง (ไม่คลายเกลียวฝาขวดสนิทและไม่เปิดฝาหม้อนึ่งขณะยังมีแรงดัน) \textbf{(CLO3, CLO4 เชื่อมโยง PLO1, PLO5, PLO6)}
\end{enumerate}

\noindent\textbf{3. กิจกรรมการจัดการเรียนรู้เชิงรุก (Active Learning \& OBE)}
\begin{itemize}
    \item \textbf{Case Study (Problem Solving)}\quad วิเคราะห์เหตุการณ์จริงในห้องแล็บจุลชีววิทยา ``เมื่อเข็มวัดความดันของเครื่อง Autoclave ชี้ที่ $15\text{ psi}$ ครบ 15 นาที แต่เมื่อนำอาหารเลี้ยงเชื้อออกมาบ่มกลับพบเชื้อปนเปื้อนเต็มขวด และเทปวัดผล (Autoclave Tape) ไม่เปลี่ยนเป็นสีดำ สาเหตุทางฟิสิกส์คืออะไร''
\end{itemize}

\noindent\textbf{4. สื่อการเรียนรู้และแหล่งเรียนรู้}
\begin{itemize}
    \item เอกสารคำสอนบทที่ 6 กราฟความดันไออิ่มตัวของน้ำ และหม้อนึ่งความดันไอ Autoclave ประจำห้องปฏิบัติการ
\end{itemize}

\noindent\textbf{5. การวัดและประเมินผลประจำบท}
\begin{itemize}
    \item การวิเคราะห์กรณีศึกษาเรื่อง Autoclave และการคำนวณกฎของแก๊ส \textbf{(CLO1, CLO2)}
    \item รายงานปฏิบัติการที่ 6 และการเตรียมความพร้อมก่อนสอบกลางภาคในสัปดาห์ที่ 7 \textbf{(CLO3, CLO4)}
\end{itemize}
\end{planbox}
\clearpage

\chapter{อุณหพลศาสตร์ ความดันไอ และหม้อนึ่งความดันไอ}
\index{อุณหพลศาสตร์}\index{ความดันไอ}\index{หม้อนึ่งความดันไอ Autoclave}\index{เอนโดสปอร์}

\section{ความดัน กฎของแก๊ส และความสัมพันธ์ระหว่างความดันกับอุณหภูมิ}

\deepskinfig{figures/sec_3d/ch06_s1.jpg}{ความสัมพันธ์ระหว่างความดันเกจ ($P_{\text{gauge}}$) ความดันบรรยากาศ ($P_{\text{atm}}$) และความดันสัมบูรณ์ ($P_{\text{abs}}$) ตามกฎของแก๊สในภาชนะปิด}

\textbf{ความดัน} (\textbf{Pressure}, $P$) คือขนาดของแรงที่กระทำตั้งฉากต่อหนึ่งหน่วยพื้นที่ ($P = F/A$) มีหน่วยเอสไอคือ \textbf{พาสคัล} ($\text{Pa}$ หรือ $\text{N/m}^2$) ในการใช้งานเครื่องมือทางจุลชีววิทยา เช่น หม้อนึ่งความดันไอ (\textbf{Autoclave}) หน้าปัดเครื่องมักแสดงหน่วยความดันเป็น\textbf{ปอนด์ต่อตารางนิ้ว} ($\text{psi}$) หรือ\textbf{บรรยากาศ} ($\text{atm}$) โดยที่ความดันบรรยากาศปกติที่ระดับน้ำทะเลมีค่าเท่ากับ

\begin{equation}
    1\text{ atm} = 101.3\text{ kPa} \approx 14.7\text{ psi} \approx 15\text{ psi}
\end{equation}

สิ่งสำคัญที่นักศึกษาต้องระวังคือ ตัวเลขที่อ่านได้จากหน้าปัดเครื่องมือวัดความดันทั่วไปคือ \textbf{ความดันเกจ} (\textbf{Gauge Pressure}, $P_{\text{gauge}}$) ซึ่งเป็นความดันส่วนที่เกินไปจากความดันบรรยากาศ ดังนั้น หากต้องการหา\textbf{ความดันสัมบูรณ์จริง} (\textbf{Absolute Pressure}, $P_{\text{abs}}$) ภายในหม้อนึ่ง จะต้องนำความดันบรรยากาศมารวมด้วยเสมอ ดังแสดงใน\mbox{ภาพที่} 6.1 และสมการที่ (6.2)

\begin{equation}
    P_{\text{abs}} = P_{\text{atm}} + P_{\text{gauge}}
\end{equation}

สำหรับภาชนะปิดที่มีปริมาตรคงที่ ความดันสัมบูรณ์ ($P$) จะแปรผันตรงกับอุณหภูมิในหน่วยเคลวิน ($T_{\text{K}}$) ตามกฎของเกย์-ลูสแซก ดังสมการที่ (6.3)

\begin{equation}
    \frac{P_1}{T_1} = \frac{P_2}{T_2} \qquad \text{(ต้องแทนค่า } T \text{ ในหน่วยเคลวิน } \text{K เสมอ)}
\end{equation}

\section{ฟิสิกส์ของหม้อนึ่งความดันไอ ทำไมต้อง 121 องศาเซลเซียส ที่ 15 psi}

\deepskinfig{figures/sec_3d/ch06_s2.jpg}{ภาพตัดขวางหม้อนึ่งความดันไอ (Autoclave) ที่สภาวะมาตรฐาน $121\,^\circ\text{C}$ ความดันเกจ $15\text{ psi}$ และการควบแน่นคายความร้อนแฝงทำลายเอนโดสปอร์ของแบคทีเรีย}

การต้มน้ำในหม้อเปิดที่ความดันบรรยากาศปกติ ($1\text{ atm}$ หรือความดันเกจ $0\text{ psi}$) น้ำจะเดือดที่อุณหภูมิสูงสุดเพียง $100\,^\circ\text{C}$ เท่านั้น ไม่ว่าจะเปิดไฟแรงเพียงใดอุณหภูมิของน้ำเดือดก็จะไม่สูงเกิน $100\,^\circ\text{C}$ ซึ่งอุณหภูมิดังกล่าวสามารถฆ่าเซลล์แบคทีเรียทั่วไป (\textbf{Vegetative Cells}) ได้ แต่\textbf{ไม่สามารถทำลายเอนโดสปอร์} (\textbf{Endospores}) ที่ทนความร้อนสูงของแบคทีเรียกลุ่ม \textit{Bacillus} และ \textit{Clostridium} ได้

เพื่อยกจุดเดือดของน้ำให้สูงขึ้นถึง $121\,^\circ\text{C}$ เราจึงต้องต้มน้ำในภาชนะปิดสนิทที่เรียกว่า \textbf{หม้อนึ่งความดันไอ} (\textbf{Autoclave}) เมื่อไอน้ำถูกกักไว้จนเกิดความดันเกจเพิ่มขึ้นเป็น $15\text{ psi}$ (ซึ่งคิดเป็นความดันสัมบูรณ์ $P_{\text{abs}} = 14.7 + 15 \approx 29.7\text{ psi}$ หรือประมาณ $2\text{ atm}$) จุดเดือดของน้ำจะขยับสูงขึ้นจาก $100\,^\circ\text{C}$ ไปอยู่ที่ $121\,^\circ\text{C}$ พอดี ดังแสดงใน\mbox{ภาพที่} 6.2 และ\mbox{ภาพที่} 6.3 และเมื่อไอน้ำอิ่มตัวที่ $121\,^\circ\text{C}$ สัมผัสกับขวดอาหารเลี้ยงเชื้อหรือเครื่องมือแพทย์ที่เย็นกว่า ไอน้ำจะควบแน่นกลายเป็นหยดน้ำพร้อมกับคาย\textbf{ความร้อนแฝงของการกลายเป็นไอ} ($L_v = 2{,}260\text{ J/g}$) ออกมาเข้าทำลายโครงสร้างโปรตีนและเอนโดสปอร์ให้ตายอย่างสมบูรณ์ภายในเวลา 15--20 นาที

\begin{figure}[H]
\centering
\resizebox{0.88\textwidth}{!}{%
\begin{tikzpicture}[font=\sffamily, >=Stealth]
    % Left: Curve of Boiling Point vs Gauge Pressure
    \begin{scope}[shift={(0,0)}]
        \draw[->, line width=1.4pt, color=charcoal] (0,0) -- (6.8,0) node[right, font=\small\bfseries] {ความดันเกจ ($\text{psi}$)};
        \draw[->, line width=1.4pt, color=charcoal] (0,0) -- (0,4.5) node[above, font=\small\bfseries] {จุดเดือดของน้ำ ($^\circ\text{C}$)};

        % Grid & Curve
        \draw[physSlate!30, dashed] (0,1.5) -- (6.2,1.5);
        \draw[physSlate!30, dashed] (3.0,0) -- (3.0,3.1) -- (0,3.1);
        \draw[physSlate!30, dashed] (5.5,0) -- (5.5,4.0) -- (0,4.0);

        \draw[physRoyalBlue, line width=2.2pt] (0,1.5) .. controls (1.5,2.4) and (3.0,3.1) .. (5.5,4.0);

        \fill[physCyan] (0,1.5) circle (3.2pt);
        \node[left, font=\footnotesize\bfseries, color=physRoyalBlue] at (0,1.5) {$100\,^\circ\text{C}$};
        \node[below, font=\footnotesize\bfseries] at (0,0) {$0\text{ psi}$};

        \fill[crimsonAlert] (3.0,3.1) circle (3.8pt);
        \node[left, font=\footnotesize\bfseries, color=crimsonAlert] at (0,3.1) {$121\,^\circ\text{C}$};
        \node[below, font=\footnotesize\bfseries, color=crimsonAlert] at (3.0,0) {$15\text{ psi}$};
        \node[above left, font=\footnotesize\bfseries, fill=boxRedBg, draw=crimsonAlert, rounded corners=2pt, inner sep=2pt] at (3.0,3.2) {จุดมาตรฐาน Autoclave};

        \fill[physAmber] (5.5,4.0) circle (3.2pt);
        \node[left, font=\footnotesize\bfseries, color=physAmber!90!black] at (0,4.0) {$134\,^\circ\text{C}$};
        \node[below, font=\footnotesize\bfseries] at (5.5,0) {$30\text{ psi}$};
    \end{scope}

    % Right: Air Pocket Problem inside Autoclave
    \begin{scope}[shift={(8.4,0)}]
        \draw[fill=boxGreenBg, draw=bioEmerald, line width=1.2pt, rounded corners=4pt] (0,2.3) rectangle (6.2,4.3);
        \node[font=\small\bfseries, color=bioEmerald!90!black] at (3.1,3.9) {(ก) ไล่อากาศออกหมด 100\%};
        \node[font=\footnotesize, align=center, color=charcoal] at (3.1,2.95) {$P_{\text{gauge}} = 15\text{ psi}$ จากไอน้ำบริสุทธิ์\\อุณหภูมิจริงถึง $121\,^\circ\text{C}$ $\rightarrow$ \textbf{สปอร์ตายหมด}};

        \draw[fill=boxRedBg, draw=crimsonAlert, line width=1.2pt, rounded corners=4pt] (0,0.0) rectangle (6.2,2.0);
        \node[font=\small\bfseries, color=crimsonAlert] at (3.1,1.6) {(ข) ไม่ได้ไล่อากาศออก (มีอากาศค้าง)};
        \node[font=\footnotesize, align=center, color=charcoal] at (3.1,0.65) {$P_{\text{gauge}} = 15\text{ psi}$ (ความดันอากาศ $+$ ไอน้ำ)\\อุณหภูมิไอน้ำแค่ $108\text{--}112\,^\circ\text{C}$ $\rightarrow$ \textbf{ฆ่าเชื้อล้มเหลว!}};
    \end{scope}
\end{tikzpicture}%
}
\caption{กราฟความสัมพันธ์ระหว่างความดันเกจกับจุดเดือดของน้ำ และผลกระทบของอากาศตกค้างภายในหม้อนึ่งความดันไอ}
\end{figure}

\vspace{4pt}
\arucomarkerbadge{5}{6}{หม้อนึ่งความดันไอ Autoclave (121°C, 15 psi) และการไล่อากาศ}{แสดงโครงสร้างภายในหม้อนึ่งความดันไอ การควบแน่นคายความร้อนแฝงของไอน้ำอิ่มตัว และผลเสียของการมีอากาศตกค้างต่อการฆ่าเชื้อสปอร์}{ความดันเกจ $P_{\text{gauge}}$ ($\text{psi}$), ร้อยละการไล่อากาศออก (\%), เวลาฆ่าเชื้อ ($\text{min}$)}


\begin{warningbox}[กฎความปลอดภัยเหล็ก 3 ข้อในการใช้หม้อนึ่งความดันไอ Autoclave]
\begin{enumerate}
    \item \textbf{ต้องไล่อากาศออกจากหม้อนึ่งให้หมดก่อนปิดวาล์ว}\quad ตามกฎความดันย่อยของดอลตัน ($P_{\text{total}} = P_{\text{steam}} + P_{\text{air}}$) หากมีอากาศปนอยู่ ความดัน $15\text{ psi}$ ที่อ่านได้จะเกิดจากอากาศผสมไอน้ำ ทำให้อุณหภูมิจริงไม่ถึง $121\,^\circ\text{C}$
    \item \textbf{ต้องคลายเกลียวฝาขวดอาหารเลี้ยงเชื้อออกประมาณครึ่งรอบเสมอ}\quad ห้ามบิดฝาขวดแก้วแน่นสนิทเข้าเครื่อง Autoclave เพราะแรงดันไอน้ำจะทำให้ขวดระเบิด
    \item \textbf{ห้ามเปิดฝาหม้อนึ่งจนกว่าเข็มความดันเกจจะลดลงมาอยู่ที่ $0\text{ psi}$}\quad และอุณหภูมิลดลงต่ำกว่า $80\,^\circ\text{C}$ มิฉะนั้นอาหารเหลวที่ร้อนจัดจะเกิดการเดือดพุ่งทะลัก (\textbf{Superheated Boiling}) ลวกผู้ปฏิบัติงาน
\end{enumerate}
\end{warningbox}

\begin{worked}[การคำนวณความดันสัมบูรณ์และการเปลี่ยนแปลงความดันตามอุณหภูมิ]
\textbf{โจทย์ปัญหา}\quad ขวดแก้วใบหนึ่งถูกปิดฝาสนิทโดยเผลอไม่ได้คลายเกลียวที่อุณหภูมิห้อง $27\,^\circ\text{C}$ (ความดันภายในขวดเริ่มต้น $P_1 = 1.0\text{ atm}$) เมื่อนำเข้าไปนึ่งในเครื่อง Autoclave จนอุณหภูมิสูงขึ้นเป็น $121\,^\circ\text{C}$ ความดันแก๊สภายในขวดจะเพิ่มขึ้นเป็นกี่บรรยากาศ ($\text{atm}$)

\textbf{วิธีทำทีละขั้นตอนฉบับเข้าใจง่าย}
\begin{enumerate}
    \item \textbf{ขั้นที่ 1 แปลงอุณหภูมิเซลเซียสให้เป็นเคลวิน ($T_{\text{K}} = T_{\text{C}} + 273$)}
    \[
        T_1 = 27 + 273 = 300\text{ K} \qquad \text{และ} \qquad T_2 = 121 + 273 = 394\text{ K}
    \]
    \item \textbf{ขั้นที่ 2 ใช้สมการสัดส่วนตรง $\dfrac{P_1}{T_1} = \dfrac{P_2}{T_2}$ เพื่อหาค่า $P_2$}
    \[
        P_2 = P_1 \times \frac{T_2}{T_1} = 1.0\text{ atm} \times \frac{394\text{ K}}{300\text{ K}} = 1.31\text{ atm}
    \]
    หากรวมกับความดันไอน้ำของของเหลวที่เดือดภายในขวด ความดันรวมภายในขวดที่ปิดสนิทจะพุ่งสูงเกิน $3\text{ atm}$ จึงเสี่ยงต่อการระเบิดอย่างยิ่ง
\end{enumerate}
\textbf{คำตอบ}\quad เฉพาะความดันของอากาศแห้งภายในขวดจะเพิ่มขึ้นเป็น $1.31\text{ atm}$ (และเมื่อรวมความดันไอน้ำจะสูงกว่า $3\text{ atm}$)
\end{worked}

\section{ปฏิบัติการประจำสัปดาห์ที่ 6}

\deepskinfig{figures/sec_3d/ch06_s3.jpg}{ชุดปฏิบัติการที่ 6 สำหรับศึกษาการเดือดของน้ำที่ความดันต่ำและการตรวจสอบประสิทธิภาพการฆ่าเชื้อด้วยเทปวัดผลความร้อน (Autoclave Indicator Tape)}

\begin{labbox}[ปฏิบัติการที่ 6\quad ความดันไอ จุดเดือด และกรณีศึกษาการทำงานของ Autoclave (3 ชั่วโมง)]
\noindent\textbf{ผลลัพธ์การเรียนรู้ที่ประเมิน}\quad CLO1 (ความดันและอุณหพลศาสตร์), CLO2 (วิเคราะห์การทำงานของ Autoclave), CLO3 (ความปลอดภัยภาชนะความดันสูง), CLO4 (รายงานกลุ่ม)

\vspace{4pt}
\noindent\textbf{1. วัสดุและอุปกรณ์ (ดังแสดงใน\mbox{ภาพที่} 6.4)}
\begin{itemize}
    \item ชุดสาธิตการเดือดของน้ำภายใต้ความดันต่ำ (กระบอกฉีดยาขนาดใหญ่หรือขวดก้นกลมต่อปั๊มสุญญากาศ) และเทอร์โมมิเตอร์
    \item หม้อนึ่งความดันไอ (Autoclave) ประจำห้องปฏิบัติการ, เทปวัดผลความร้อน (Autoclave Indicator Tape) และขวดใส่สารละลายตัวอย่าง
\end{itemize}

\noindent\textbf{2. ขั้นตอนการทดลอง}
\begin{enumerate}
    \item \textbf{ตอนที่ 1 การเดือดของน้ำที่ความดันต่ำ}\quad ดูดน้ำอุ่นอุณหภูมิประมาณ $70\,^\circ\text{C}$ เข้ากระบอกฉีดยาใส อุดปลายกระบอกแล้วดึงก้านสูบออกเพื่อลดความดันภายใน สังเกตฟองการเดือดของน้ำที่อุณหภูมิต่ำกว่า $100\,^\circ\text{C}$
    \item \textbf{ตอนที่ 2 การบันทึกกราฟรอบการทำงานจริงของเครื่อง Autoclave}\quad เตรียมขวดน้ำกลั่นที่คลายเกลียวฝาครึ่งรอบและติดเทป Autoclave Tape นำเข้าเครื่อง Autoclave บันทึกค่าความดันเกจ ($\text{psi}$) และอุณหภูมิ ($^\circ\text{C}$) ทุกๆ 5 นาที ตั้งแต่ช่วงไล่อากาศ ช่วงเพิ่มความดัน ช่วงคงที่ที่ $121\,^\circ\text{C}$ ($15\text{ psi}$) นาน 15 นาที จนถึงช่วงลดความดัน
\end{enumerate}

\noindent\textbf{3. ตารางบันทึกผลการทดลอง}
\begin{center}
\footnotesize\setlength{\tabcolsep}{3.5pt}
\begin{tabularx}{\linewidth}{>{\raggedright\arraybackslash}X c c c}
\toprule
\textbf{ขั้นตอนการทำงานของ Autoclave} & \textbf{เวลาที่ผ่านไป ($\text{min}$)} & \textbf{ความดันเกจ ($\text{psi}$)} & \textbf{อุณหภูมิที่อ่านได้ ($^\circ\text{C}$)} \\
\midrule
1. เริ่มต้นเปิดเครื่องและไล่อากาศ & $0\text{ min}$ & ............ & ............ \\
2. ปิดวาล์วไล่อากาศและเพิ่มความดัน & $10\text{ min}$ & ............ & ............ \\
3. เข้าสู่สภาวะฆ่าเชื้อมาตรฐาน & $15\text{ min}$ & ............ & ............ \\
4. ครบเวลาฆ่าเชื้อ 15 นาที & $30\text{ min}$ & ............ & ............ \\
5. ความดันลดลงปลอดภัยพร้อมเปิดฝา & $45\text{ min}$ & ............ & ............ \\
\bottomrule
\end{tabularx}
\end{center}
\end{labbox}

\section*{คำถามท้ายบทที่ 6 และแนวข้อสอบกลางภาค (สัปดาห์ที่ 7)}
\begin{enumerate}
    \item เมื่ออ่านหน้าปัดความดันของเครื่อง Autoclave ได้ค่าความดันเกจเท่ากับ $15\text{ psi}$ จงหาว่าความดันสัมบูรณ์ภายในหม้อนึ่งมีค่ากี่ $\text{psi}$ และคิดเป็นกี่บรรยากาศ ($\text{atm}$)
    \item เพราะเหตุใดไอน้ำอิ่มตัวที่อุณหภูมิ $121\,^\circ\text{C}$ ในเครื่อง Autoclave จึงสามารถฆ่าเอนโดสปอร์ได้ภายในเวลา 15 นาที ในขณะที่ตู้อบลมร้อนแห้ง (Hot Air Oven) ต้องใช้อุณหภูมิสูงถึง $160\text{--}180\,^\circ\text{C}$ และใช้เวลานานถึง $1\text{--}2$ ชั่วโมง จงอธิบายด้วยหลักการความร้อนแฝงของการควบแน่น
\end{enumerate}
